\usepackage{booktabs}
\usepackage{longtable}
\usepackage{array}
\usepackage{tabularx}
\usepackage{caption}
\usepackage{fvextra}
\usepackage{enumitem}
\usepackage{fancyhdr}
\usepackage{xcolor}
\usepackage{hyperref}
\usepackage{titlesec}
\usepackage{setspace}
\usepackage{float}
\usepackage{graphicx}


% 行内代码样式：用于 Markdown 反引号 `code`。
% 不使用 \colorbox：单个颜色盒不能跨行，会让长行内代码越过页边界。
% filters/inline-code-breaks.lua 会在标点与固定长度处插入不可见断点。
\definecolor{InlineCodeText}{HTML}{6F42C1}
\DeclareRobustCommand{\texttt}[1]{%
  \begingroup
  \textcolor{InlineCodeText}{\ttfamily\bfseries #1}%
  \endgroup
}

\DefineVerbatimEnvironment{Highlighting}{Verbatim}{breaklines,breakanywhere,commandchars=\\\{\}}
\fvset{fontsize=\small,breaklines=true,breakanywhere=true}
\setlist{nosep}
\setlength{\parindent}{2em}
\setlength{\parskip}{0.35em}
\setlength{\emergencystretch}{2em}
\renewcommand{\arraystretch}{1.2}
\captionsetup{font=small}

\pagestyle{fancy}
\fancyhf{}
\fancyhead[L]{Banewfn Manual}
\fancyhead[R]{\nouppercase{\leftmark}}
\fancyfoot[C]{\thepage}
\setlength{\headheight}{14pt}

\titleformat{\section}{\Large\bfseries}{\thesection}{0.8em}{}
\titleformat{\subsection}{\large\bfseries}{\thesubsection}{0.8em}{}
\titleformat{\subsubsection}{\normalsize\bfseries}{\thesubsubsection}{0.8em}{}
